% Computational discovery chapter: continued fractions and Euler sums
\chapter{Continued Fractions and Euler Sums}
\label{ch:discovery_continued_fractions}

\section{Rational Approximation}

Repeatedly taking the integer part of a number and then inverting the remainder
produces a continued fraction. For the rational approximation $355/113$ to $\pi$:

\begin{lstlisting}[language=Python]
from fractions import Fraction
from euler import continued_fraction, continued_fraction_convergents

a = continued_fraction(Fraction(355, 113))
print(a)
print(continued_fraction_convergents(a))
\end{lstlisting}

The convergents $3$, $22/7$, and $355/113$ appear naturally. This is an example of
computation revealing that good approximations are not accidental.

Euler's continued fraction for $e$ has the pattern

\[
e=[2;1,2,1,1,4,1,1,6,1,1,8,\ldots].
\]

\begin{lstlisting}[language=Python]
from fractions import Fraction
from euler import continued_fraction_convergents

coefficients = [2, 1, 2, 1, 1, 4, 1, 1, 6, 1, 1, 8]
for convergent in continued_fraction_convergents(coefficients):
    print(convergent, float(convergent))
\end{lstlisting}

\section{Euler Sums}

The harmonic numbers $H_n$ occur throughout Euler's work. A computer can suggest
the non-obvious identity

\[
\sum_{n=1}^{\infty}\frac{H_n}{n^2}=2\zeta(3).
\]

\begin{lstlisting}[language=Python]
from euler import euler_sum_harmonic_over_square

for n in (100, 1000, 10000):
    value = euler_sum_harmonic_over_square(n)
    print(n, value, value - 2.4041138063191885)
\end{lstlisting}

The experiment is evidence, not proof. A proof can reorganize the double sum:

\[
\sum_{n\geq1}\frac{H_n}{n^2}
=\sum_{n\geq1}\sum_{k=1}^{n}\frac{1}{k n^2},
\]

then use symmetry and telescoping identities. This is precisely the kind of
rearrangement through which Euler turned numerical patterns into analysis.

\section{Exercises}
\begin{exercise}
Find the first ten convergents of $\sqrt2=[1;2,2,2,\ldots]$ and plot their errors.
\end{exercise}
\begin{exercise}
Compare the approximation to $\pi$ from $22/7$ and $355/113$ with decimal truncation.
\end{exercise}
\begin{exercise}
Estimate the tail of the Euler sum and explain why direct summation is slow.
\end{exercise}
